% SHARD-ID: AQM-21-GRASSMANN-WEYL-FERMION-DERIVATIONS
% SHARD-TITLE: Grassmann-Weyl Fermion Displacement Derivations
% SHARD-SUMMARY: Proves the unitarity of Cahill--Glauber fermion displacement operators.
% SHARD-SUMMARY: Derives the CAR displacement identities and Grassmann-valued ray product.
% SHARD-SUMMARY: Records the boundary between Grassmann-extended CAR unitarity and ordinary Hilbert-space unitarity.
% SHARD-KEYWORDS: Grassmann, CAR, fermion displacement, Baker-Hausdorff, coherent states, unitary

\section{Grassmann-Weyl Fermion Displacement Derivations}
\label{sec:grassmann-weyl-fermion-displacement-derivations}

\subsection{Status and Dependencies}

This shard proves the Lamport items \(L2\)--\(T2\) stated in
Section~\ref{sec:grassmann-weyl-fermion-displacements}.  It uses the same
Cahill--Glauber source anchors from
\path{references/canonical_fields/cahill_glauber_physics_9808029v1_source/fxxx.tex}:
lines 367--390 for \(D(\gamma)\), lines 391--414 for the displacement
identities, lines 416--464 for the Baker--Hausdorff and ray-composition
formula, and lines 835--854 for the ordinary-operator caveat.

\[
\begin{array}{rcl}
D1\text{--}D7&:&\text{definitions from Section }
  \ref{sec:grassmann-weyl-fermion-displacements}.\\
L1&:&\sigma\text{ is the central Grassmann two-cocycle form.}\\
L2&:&K(\gamma)^\dagger=-K(\gamma),\text{ so }D(\gamma)\text{ is unitary.}\\
L3&:&D(\gamma)^\dagger a_nD(\gamma)=a_n+\gamma_n
      \text{ and similarly for }a_n^\dagger.\\
L4&:&[K(\alpha),K(\beta)]=\sigma(\alpha,\beta).\\
T1&:&D(\gamma)\text{ represents the abstract displacement ray algebra.}\\
T2&:&\text{The theorem is Grassmann-extended CAR-unitary, not ordinary
      scalar-projective Hilbert-unitary.}
\end{array}
\]

\subsection{Unitarity L2}

\paragraph{Lemma \(L2\).}
\[
  K(\gamma)^\dagger=-K(\gamma),
  \qquad
  D(\gamma)^\dagger D(\gamma)=D(\gamma)D(\gamma)^\dagger=1 .
\]

\emph{Proof.}
Using the adjoint convention from \(D3\),
\[
  (a_i^\dagger\gamma_i)^\dagger=\gamma_i^*a_i,\qquad
  (\gamma_i^*a_i)^\dagger=a_i^\dagger\gamma_i .
\]
Therefore
\[
  K(\gamma)^\dagger
  =
  \sum_i(\gamma_i^*a_i-a_i^\dagger\gamma_i)
  =
  -K(\gamma).
\]
Hence \(D(\gamma)^\dagger=e^{-K(\gamma)}=D(-\gamma)\), and multiplying
\(e^{-K(\gamma)}e^{K(\gamma)}\) or \(e^{K(\gamma)}e^{-K(\gamma)}\) gives \(1\).

\subsection{Displacement of the CAR Generators L3}

\paragraph{Lemma \(L3\).}
For every \(n\in I\),
\[
  D(\gamma)^\dagger a_nD(\gamma)=a_n+\gamma_n,\qquad
  D(\gamma)^\dagger a_n^\dagger D(\gamma)=a_n^\dagger+\gamma_n^* .
\]

\emph{Proof.}
The conjugation identity \(e^{-K}Be^K=B+[B,K]+\frac12[[B,K],K]+\cdots\)
applies because \(K\) is even.  For \(i\neq n\), the \(i\)-summand of
\(K(\gamma)\) commutes with \(a_n\): the sign from moving \(a_n\) through the
CAR generator is cancelled by the sign from moving it through the Grassmann
coefficient.  For \(i=n\),
\[
\begin{aligned}
 [a_n,a_n^\dagger\gamma_n-\gamma_n^*a_n]
 &=
 a_na_n^\dagger\gamma_n-a_n^\dagger\gamma_na_n\\
 &=
 (1-a_n^\dagger a_n)\gamma_n-a_n^\dagger\gamma_na_n\\
 &=
 \gamma_n-a_n^\dagger a_n\gamma_n+a_n^\dagger a_n\gamma_n\\
 &=\gamma_n .
\end{aligned}
\]
The element \(\gamma_n\) commutes with \(K(\gamma)\), so all higher
commutators vanish.  Thus \(D(\gamma)^\dagger a_nD(\gamma)=a_n+\gamma_n\).
The same calculation with \(a_n^\dagger\) gives
\[
\begin{aligned}
 [a_n^\dagger,a_n^\dagger\gamma_n-\gamma_n^*a_n]
 &=
 -a_n^\dagger\gamma_n^*a_n+\gamma_n^*a_na_n^\dagger\\
 &=
 \gamma_n^*a_n^\dagger a_n+\gamma_n^*(1-a_n^\dagger a_n)\\
 &=\gamma_n^* ,
\end{aligned}
\]
and again higher commutators vanish.

\subsection{The Grassmann Commutator L4}

\paragraph{Lemma \(L4\).}
For all parameters \(\alpha,\beta\),
\[
  [K(\alpha),K(\beta)]=\sigma(\alpha,\beta).
\]

\emph{Proof.}
It is enough to compute one mode, since different modes commute after the same
double sign cancellation used in Lemma \(L3\).  Put
\[
  p=a^\dagger\alpha,\quad q=\alpha^*a,\quad
  r=a^\dagger\beta,\quad s=\beta^*a .
\]
Then \(K(\alpha)=p-q\) and \(K(\beta)=r-s\).  The products \(pr,rp,qs,sq\)
vanish because \((a^\dagger)^2=a^2=0\).  The mixed commutators are
\[
\begin{aligned}
 [p,s]
 &=a^\dagger\alpha\beta^*a-\beta^*aa^\dagger\alpha\\
 &=-\beta^*\alpha\,a^\dagger a-\beta^*(1-a^\dagger a)\alpha\\
 &=-\beta^*\alpha,
\end{aligned}
\]
and
\[
\begin{aligned}
 [q,r]
 &=\alpha^*aa^\dagger\beta-a^\dagger\beta\alpha^*a\\
 &=\alpha^*(1-a^\dagger a)\beta+\alpha^*\beta\,a^\dagger a\\
 &=\alpha^*\beta .
\end{aligned}
\]
Therefore
\[
  [K(\alpha),K(\beta)]
  =
  [p-q,r-s]
  =
  -[p,s]-[q,r]
  =
  \beta^*\alpha-\alpha^*\beta .
\]
Summing over \(i\in I\) gives the displayed formula.

\subsection{Representation Theorem T1}

\paragraph{Theorem \(T1\).}
The assignment
\[
  \Pi(\mathsf D(\gamma))=
  \exp\!\left(\sum_i(a_i^\dagger\gamma_i-\gamma_i^*a_i)\right)
\]
is a unitary representation of the abstract Grassmann displacement ray algebra
in the Grassmann-extended finite CAR algebra.

\emph{Proof.}
Lemma \(L2\) proves unitarity and the \(*\)-relation
\(D(\gamma)^\dagger=D(-\gamma)\).  Lemma \(L4\) says that the commutator
\([K(\alpha),K(\beta)]\) is the central element \(\sigma(\alpha,\beta)\).
The Baker--Hausdorff identity for central commutator gives
\[
\begin{aligned}
  e^{K(\alpha)}e^{K(\beta)}
  &=
  e^{K(\alpha)+K(\beta)}
  e^{\frac12[K(\alpha),K(\beta)]}\\
  &=
  e^{K(\alpha+\beta)}
  e^{\frac12\sigma(\alpha,\beta)}\\
  &=
  D(\alpha+\beta)c(\alpha,\beta).
\end{aligned}
\]
This is exactly the defining product relation of \(\mathsf D_I\).  Lemma
\(L1\) supplies associativity of the abstract relation, so \(\Pi\) is a
well-defined representation.

\subsection{Boundary of the Statement T2}

\paragraph{Theorem \(T2\).}
The positive theorem \(T1\) is not an ordinary projective representation
\(\gamma\mapsto U_\gamma\in U(\mathcal H)\) with complex \(U(1)\) multiplier.
It is a unitary representation in \(\mathcal A_I\), the CAR algebra extended by
Grassmann coefficients.

\emph{Proof.}
The multiplier in \(T1\) is
\[
  c(\alpha,\beta)=
  \exp\!\left(\frac12\sum_i(\beta_i^*\alpha_i-\alpha_i^*\beta_i)\right),
\]
which is an even Grassmann unit.  Cahill--Glauber's notation records that no
nonzero Grassmann monomial is an ordinary complex number, and it later
separates physical quantum operators from Grassmann variables.  Therefore
unless the Grassmann coefficient algebra is included in the target, the
multiplier is not a complex phase.

More explicitly, if \(\rho:\Lambda_I\to B(\mathcal H)\) is a Hilbert-space
\(*\)-representation of the Grassmann coefficient algebra, then
\(\rho(\gamma_i)=0\) for each \(i\).  Indeed, with \(b=\rho(\gamma_i)\),
\[
  b^2=0,\qquad bb^\dagger+b^\dagger b=0,
\]
and hence for every \(v\in\mathcal H\),
\[
  0=\langle v,(bb^\dagger+b^\dagger b)v\rangle
   =\|b^\dagger v\|^2+\|bv\|^2,
\]
so \(b=0\).  Thus the ordinary complex Hilbert representation collapses the
Grassmann cocycle.  The representation found in \(T1\) avoids this collapse by
keeping the Grassmann variables as formal odd coefficients anticommuting with
the concrete CAR generators.
